\section{Results and Discussion}
\label{sec:results}

We report the experiments in three steps. We first audit where the FAR-Trans baselines sit on the three axes and across the four declared bands. We then evaluate risk-based segmentation and its two exclusive levers on the aggregate balance, the per-band suitability, and the per-split consistency. We close by explaining why correcting the data beats penalising the model. Aggregate metrics over the 69 splits are in Table~\ref{tab:aggregate}, per-band suitability in Table~\ref{tab:perband}.

\begin{table*}[!t]
\centering
\small
\begin{tabular}{lrrrrr}
\toprule
Configuration & nDCG@10 & ROI@10 & PC@10 & Balance & PC-lift@10 \\
\midrule
Random Forest baseline                          & $0.019$ & $+0.014$ & $0.667$ & $0.054$ & $1.06\times$ \\
LightGCN baseline                               & $0.330$ & $-0.005$ & $0.794$ & $0.428$ & $1.17\times$ \\
\midrule
Segmentation                                    & $0.332$ & $-0.007$ & $0.821$ & $0.411$ & $1.19\times$ \\
Segmentation $+$ margin (model lever)           & $0.309$ & $-0.007$ & $0.918$ & $0.407$ & $1.40\times$ \\
Segmentation $+$ regrouping (data lever)        & $0.352$ & $-0.006$ & $0.949$ & $\mathbf{0.452}$ & n/a \\
Segmentation $+$ regrouping $+$ margin (combined) & $0.333$ & $-0.005$ & $0.982$ & $0.444$ & n/a \\
\bottomrule
\end{tabular}
\caption{Aggregate metrics averaged over the 69 temporal splits. ROI@10 is per month. Balance is the harmonic mean of the three axes with ROI rescaled (Section~\ref{sec:exp:metrics}). PC-lift@10 $=$ PC@10$/\pi(b)$ and is n/a for the regrouped configurations, whose shifted band mix makes aggregate lift incomparable. Regrouping is the only intervention whose balance exceeds the LightGCN baseline ($0.452$ vs $0.428$).}
\label{tab:aggregate}
\end{table*}

\begin{table*}[!t]
\centering
\small
\begin{tabular}{lrrrrrrr}
\toprule
Band & $\pi(b)$ & RF & LightGCN & Segment. & $+$Margin & $+$Regroup & Combined \\
\midrule
Conservative & $0.65$ & $0.52$ & $0.75$ & $0.91$ & $1.19$ & $1.21$ & $1.25$ \\
Income       & $0.78$ & $0.70$ & $1.13$ & $1.20$ & $1.26$ & $1.25$ & $1.27$ \\
Balanced     & $0.76$ & $1.03$ & $1.17$ & $1.22$ & $1.28$ & $1.24$ & $1.30$ \\
Aggressive   & $0.35$ & $1.88$ & $1.35$ & $1.17$ & $1.97$ & $2.24$ & $2.68$ \\
\bottomrule
\end{tabular}
\caption{Per-band PC-lift@10 (PC@10$/\pi(b)$) for the baselines and the segmentation configurations. Values $> 1$ beat the band-conditional random baseline $\pi(b)$. Both baselines under-serve the Conservative band, and regrouping is the first configuration to clear the random baseline in every band.}
\label{tab:perband}
\end{table*}

\subsection{Baseline Audit}
\label{sec:results:baseline}
The two FAR-Trans baselines occupy opposite corners of the return-accuracy tradeoff and neither balances the three axes (Table~\ref{tab:aggregate}). Random Forest wins return ($+1.4$\%/mo) but ranks near random (nDCG@10 of $0.019$), so its balance collapses to $0.054$. LightGCN wins ranking accuracy (nDCG@10 of $0.330$, PC@10 of $0.794$) but earns a slightly negative ROI@10, for a balance of $0.428$. Aggregate PC@10 looks adequate for both (LightGCN $1.17\times$ lift, Random Forest $1.06\times$), but that single number averages over bands whose random baselines $\pi(b)$ range from $0.35$ to $0.78$ and so hides which bands are actually served.

Table~\ref{tab:perband} exposes the gap. For both baselines, per-band lift rises monotonically with the declared band: $0.52 \to 0.70 \to 1.03 \to 1.88$ for Random Forest and $0.75 \to 1.13 \to 1.17 \to 1.35$ for LightGCN. Recommendations skew towards higher-risk assets, so both baselines under-serve the Conservative band (below the random baseline), and Random Forest also under-serves Income. Aggressive's high lift is partly mechanical, since $\pi(\text{Aggressive}) = 0.35$ is the lowest bar to clear. Neither baseline carries band labels or a band-aware objective, so per-band coherence is a passive byproduct of training rather than a controlled output. This is the deficit a band-aware approach must close.

\subsection{Aggregate Results of Risk-Based Segmentation}
\label{sec:results:aggregate}
Reading the ablation down Table~\ref{tab:aggregate} separates a cheap architectural step from the two levers. Segmentation alone lifts PC@10 to $0.821$ at roughly flat ranking (nDCG@10 of $0.332$) and a small ROI cost, for a balance of $0.411$, marginally below the baseline. The two levers then diverge sharply. The model-side margin loss pushes PC@10 to $0.918$ but spends accuracy to do so, dropping nDCG@10 to $0.309$. Its balance of $0.407$ sits below the baseline, so it buys coherence at the cost of the overall balance. The data-side regrouping instead lifts PC@10 to $0.949$ while \emph{also} raising nDCG@10 to $0.352$ and leaving ROI@10 essentially unchanged. Its balance of $0.452$ is the highest of any configuration and the only one above the LightGCN baseline. Combining both levers reaches the highest raw PC@10 ($0.982$) but its balance ($0.444$) falls back below regrouping alone, because the margin loss spends accuracy even on the corrected bands. The headline is therefore that of the two exclusive levers, only regrouping raises the harmonic-mean balance.

\subsection{Per-Band Suitability}
\label{sec:results:perband}
The per-band lifts in Table~\ref{tab:perband} trace where each lever spends its effort. The margin loss is largely spent rescuing the two extremes that the baselines failed: it drags Conservative from $0.75$ past the random baseline to $1.19$ and lifts Aggressive from $1.35$ to $1.97$. Segmentation alone moves Conservative only to $0.91$ (still below random) and even lowers Aggressive to $1.17$, likely because the smaller per-band training graph yields weaker embeddings in isolation. Regrouping clears the random baseline in every band ($1.21$, $1.25$, $1.24$, $2.24$), and adding the margin on top of it reaches the highest per-band lifts of all ($1.25$, $1.27$, $1.30$, $2.68$). The Conservative deficit that both baselines exhibited is thus closed by either lever, but only regrouping closes it without the accuracy penalty seen in the aggregate.

\subsection{Consistency Across Splits}
\label{sec:results:splits}
The aggregate gains are not artefacts of averaging over a few strong splits. Figure~\ref{fig:paired-deltas} shows the per-split deltas against the baseline: PC@10 is positive on every one of the 69 splits for all three configurations, and the balance delta is positive for regrouping ($+0.019$ on average) but negative for both segmentation alone and the margin loss. Figure~\ref{fig:win-rates} reports the win rate, the fraction of splits on which a configuration beats the baseline. Regrouping wins the balance on $71\%$ of splits and nDCG@10 on all of them, whereas the margin loss wins the balance on only $26\%$ of splits and nDCG@10 on a mere $7\%$: its accuracy cost is a uniform downward shift, not a handful of bad splits. The per-split picture confirms the aggregate ranking of the levers.

\begin{figure}[!t]
    \centering
    \includegraphics[width=\linewidth]{figures/paired_deltas.pdf}
    \caption{Paired per-split deltas against the LightGCN baseline for segmentation, the model lever, and the data lever. PC@10 is positive on every split, but only regrouping (data lever) yields a positive Balance delta.}
    \label{fig:paired-deltas}
\end{figure}

\begin{figure}[!t]
    \centering
    \includegraphics[width=\linewidth]{figures/paired_win_rates.pdf}
    \caption{Per-split win rates against the LightGCN baseline. Regrouping wins the Balance on $71\%$ of splits and nDCG@10 on all of them, while the margin loss wins the Balance on $26\%$ and nDCG@10 on only $7\%$.}
    \label{fig:win-rates}
\end{figure}

\subsection{Why Regrouping Beats the Margin Loss}
\label{sec:results:why}
The two levers attack the same problem from opposite ends, and the balance shows the difference. The margin loss fights the noisy declared labels during training by penalising discordant scores. It raises coherence but spends ranking accuracy doing so, and the per-band table shows that cost is largely incurred rescuing a single band. Regrouping instead corrects the labels at their source, routing and training each customer on the band implied by their revealed purchases, so coherence rises without any fight against the data. The decisive evidence is that nDCG@10 and ROI@10, both independent of the risk band, improve under regrouping: a model cannot rank held-out purchases better by relabelling bands unless the original labels were injecting noise into training. Their simultaneous improvement is direct evidence that the declared bands, not the model, were the binding constraint. The recommended operating point is therefore the data lever alone: once the bands are corrected, the margin loss is redundant, lowering the balance from $0.452$ to $0.444$ while adding only raw coherence the application does not need.
